% исправления 27.05.2005
% исправления 05.10.2007
\section{Неполные и неразрешимые теории}
        \label{noncomplete-undecidable}

Предыдущий раздел мог создать впечатление, что наугад взятая
теория скорее всего окажется полной, разрешимой, а возможно, и
конечно аксиоматизируемой. Это совсем не так.

Откуда вообще берутся в математике аксиоматические теории?
Иногда мы пытаемся построить аксиоматически теорию какой\д то
конкретной структуры (скажем, теорию действительных чисел со
сложением и умножением). В других случаях мы стараемся выделить
общие свойства различных структур. Например, аксиомы группы
фиксируют общие свойства различных групп, и с самого начала
ясно, что такая теория не должна и не может быть полной. То же
самое можно сказать и про теорию линейно упорядоченных
множеств\т полнота такой теории означала бы, что все линейно
упорядоченные множества (или группы) элементарно эквивалентны,
то есть обладают одними и теми же свойствами, выражаемыми
формулами. Это, конечно, не так.

Что касается конкретных структур, то и для них естественные
теории не всегда оказываются полными. Классический пример\т
натуральные числа со сложением и умножением. Для них имеется
естественная формальная теория (называемая формальной
арифметикой). Её аксиомы включают в себя обычные свойства
сложения и умножения, а также аксиомы индукции. Опыт показывает,
что любое рассуждение теории чисел, в котором речь идёт только о
конечных объектах, может быть формально записано в виде вывода
из аксиом этой теории. Более того, многие доказательства,
использующие бесконечные объекты (скажем, важнейшую в теории
чисел~$\zeta$\д функцию Римана),%
\index{$\zeta$\д функция Римана}%
\glossary{\Риман}
могут быть модифицированы
и погружены в эту формальную теорию. Тем не менее эта теория
неполна (и не может быть полна, как мы увидим в этом разделе).

Среди естественных неполных теорий бывают разрешимые и
        \index{Теория!линейно упорядоченных множеств}%
        \index{Теория!абелевых групп}%
        \index{Теория!групп}%
не\-раз\-ре\-ши\-мые. Например, теория линейно упорядоченных множеств
разрешима, теория абелевых групп разрешима, а теория групп
не\-раз\-ре\-ши\-ма. Подробный рассказ об этом далеко выходит за рамки
нашей книжки; написанный М.\,О.\,Рабином%
\glossary{\Рабин}
обзор соответствующих
результатов можно найти в справочнике по математической логике
(часть III, \emph{Теория рекурсии}~\cite{handbook-recursion},
глава~4).

Мы же ограничимся тремя примерами (теория равенства, теория
полугрупп, формальная арифметика).

\subsection*{Теория равенства}
   \index{Теория!равенства}%

Рассмотрим сигнатуру, содержащую единственный двуместный
предикат равенства, и теорию, состоящую из трёх аксиом равенства
      \index{Аксиомы!равенства}%
      \index{Равенство}%
(рефлексивность, симметричность и транзитивность). Эти аксиомы
рассматривались в разделе~\ref{equality-axioms}; заметим, что у
нас нет других предикатных и функциональных символов (и
связанных с ними аксиом равенства).

Моделями этой теории являются всевозможные множества с
отношениями эквивалентности. Нормальными моделями этой теории
являются множества различной мощности; поскольку никакой
дополнительной структуры нет, такая модель определяется
мощностью (с точностью до изоморфизма). Теоремами этой теории
будут формулы с равенством, истинные в множествах любой
мощности.

\begin{theorem}
        \label{equality-decidable}%
Множество теорем теории равенства яв\-ля\-ет\-ся разрешимым.
        %
\end{theorem}

\begin{proof}
        %
Заметим, что истинность формулы в нормальной модели может
зависеть от её мощности. Например, формула $\exists x\exists y
\lnot (x=y)$ ложна в одноэлементной модели и истинна во всех
остальных. Поэтому процедура элиминации кванторов в чистом виде
(без расширения сигнатуры) здесь неприменима.

Но идея остаётся той же. От чего зависит истинность формулы
этой сигнатуры (с параметрами)? Во\д первых, от значений параметров
(важно, какие параметры равны друг другу, а какие нет).
Во\д вторых, от числа элементов модели. (Если бы этой зависимости
не было, то можно было бы написать бескванторную формулу,
эквивалентную данной во всех моделях теории.)

Например, формула $\exists z\,
({\lnot(z=x)}\hm\land{\lnot(z=y)})$ при $x=y$ истинна во всех
моделях, начиная с двухэлементной, а при $x\ne y$ истинна во
всех моделях, начиная с трёхэлементной. Можно ожидать, что
модели с большим числом элементов неотличимы друг от друга и от
бесконечных моделей.

\textsf{Лемма}. Истинность формулы языка с равенством, содержащей
$k$ параметров и имеющей кванторную глубину $l$, определяется
тем, какие из параметров равны друг другу, а также мощностью
носителя, при этом все мощности, большие $k+l$, одинаковы.

Докажем лемму с помощью индукции по построению формулы.
Для атомарной (и вообще для любой бескванторной) формулы
мощность вообще не играет роли. Если утверждение леммы верно для
формул~$\varphi$ и~$\psi$, то оно очевидным образом верно и для
$\varphi\hm\land\psi$, $\varphi\hm\lor\psi$, $\varphi\hm\to\psi$
и~$\lnot\varphi$. При этом используется такой факт:
кванторная глубина (число вложенных кванторов) и число
параметров у части формулы не больше, чем у всей формулы.

Содержательный случай\т когда формула начинается с
квантора. Когда, скажем, формула
        $$
\exists x \varphi(x,x_1,\dots,x_k)
        $$
с параметрами $x_1,\dots,x_k$ будет истинной (в данной
интерпретации при данных значениях параметров)? Достаточно
попробовать в качестве~$x$ значения $x_1,\dots,x_k$ а также
какой\д нибудь элемент, отличный от всех этих значений. (Все
такие элементы ничем не отличаются.) Истинность
формулы~$\varphi(x,x_1,\dots,x_k)$ при~$x=x_i$ определяется
соотношениями между
параметрами и мощностью модели (по предположению
индукции; заметим, что число параметров увеличилось на~$1$, а
кванторная глубина уменьшилась на~$1$, так что сумма осталась
прежней). Существование элемента, отличного от всех~$x_i$,
определяется мощностью модели и числом различных элементов среди
$x_1,\dots,x_k$ (то есть в конечном счёте равенствами вида
$x_i\hm=x_j$). При этом модели всех мощностей, начиная с $k+1$,
ведут себя одинаково. Кроме того, истинность
формулы~$\varphi(x,x_1,\dots,x_k)$ при
$x\hm\notin\{x_1,\dots,x_k\}$ по предположению индукции также
определяется равенствами вида $x_i\hm=x_j$ и мощностью модели.

Квантор всеобщности рассматривается точно так же (а можно его
выразить через квантор существования и вообще не рассматривать).
Лемма доказана.

Доказательство леммы конструктивно, то есть указывает способ
узнать, будет ли формула истинной, если известно, какие её параметры
равны и какова мощность носителя. В частности, для замкнутых формул
получаем способ проверять их истинность для всех значений мощности,
то есть выводимость в теории равенства.
        %
\end{proof}

\begin{problem}
        \label{unary-equality}
Рассмотрим теорию, в сигнатуре которой есть равенство и конечное
число одноместных предикатных символов, а аксиомами являются
аксиомы равенства (включая устойчивость предикатов относительно
равенства, как в разделе~\ref{equality-axioms}).
Покажите, что эта теория разрешима.

(Указание. Можно провести элиминацию кванторов, добавив к
сигнатуре счётное число нульместных предикатных символов помимо
уже имеющихся одноместных предикатных символов $P_1,\ldots,P_n$.
А именно, для каждого $n$-битового слова $z$ и для каждого
натурального $k$ мы добавляем утверждение о том, что существует
ровно $k$ элементов $x$, для которых $(P_1(x),\ldots,P_n(x))$
равно~$z$.)
        %
\end{problem}

Эта задача показывает, что добавление одноместных предикатов в
сигнатуру не делает теорию равенства неразрешимой. Отметим, что
расширение сигнатуры (без изменения множества аксиом) может
превратить разрешимую теорию в неразрешимую: например, добавив
конечное число одноместных функциональных символов к теории
равенства, получим неразрешимую теорию (как мы вскоре увидим,
доказывая теорему Чёрча с помощью проблемы тождества для
полугрупп, с.~\pageref{church-theorem}). Добавление одного
двуместного предикатного символа также даёт неразрешимую теорию.

\subsection*{Теория полугрупп}

Наш второй пример\т теория полугрупп.%
     \index{Теория!полугрупп|defin}
Её сигнатура состоит из
равенства и единственного двуместного функционального символа,
называемого умножением; результат умножения $x$ и~$y$ мы будем
обозначать~$(xy)$.

Теория состоит из аксиом равенства (в них входит корректность
умножения: $(x_1\hm=x_2)\hm\land (y_1\hm=y_2) \hm\to
(x_1y_1\hm=x_2y_2)$; мы опускаем внешние кванторы всеобщности) и
аксиомы ассоциативности
        $$
\forall x \forall y \forall z \,  ((xy)z=x(yz)).
        $$

Нормальные модели этой теории называются \emph{полугруппами}.%
     \index{Полугруппа|defin}

\begin{theorem}
        %
Множество теорем теории полугрупп (то есть множество замкнутых формул
указанной сигнатуры, истинных во всех полугруппах) неразрешимо.
        %
\end{theorem}


\begin{proof}
        %
Нам понадобится конкретный способ задания полугрупп с помощью
образующих и соотношений. Пусть фиксировано некоторое конечное
множество, называемое \emph{алфавитом}.%
\index{Алфавит|defin}
Элементы его называют
\emph{буквами},%
\index{Буква алфавита|defin}
а конечные последовательности букв\т
\emph{словами}%
\index{Слово алфавита|defin}
(данного алфавита). На словах определена операция
соединения (приписывания), относительно которой они образуют
полугруппу, которая называется \emph{свободной полугруппой}.%
\index{Свободная полугруппа|defin}%
\index{Полугруппа!свободная|defin}
Эта полугруппа имеет нейтральный элемент\т пустое слово,
приписывание которого к любому слову не меняет последнего.

Пусть фиксирован алфавит $A$, а также конечное число
пар слов $(X_1, Y_1),\dots,(X_n,Y_n)$
этого алфавита. Два слова алфавита~$A$ назовём
эквивалентными, если одно можно превратить в другое, многократно
делая замены подслов вида $X_i\hm\leftrightarrow Y_i$. Легко
проверить, что получается отношение эквивалентности и что
операция приписывания корректно определена на классах
эквивалентности и ассоциативна. Получается полугруппа. Её
называют полугруппой с \emph{образующими} из~$A$ и
\emph{соотношениями}~$X_i=Y_i$.%
\index{Полугруппа!с образующими и соотношениями|defin}%
\index{Образующая полугруппы|defin}%
\index{Соотношение полугруппы|defin}

        \begin{problem}
Сколько элементов в полугруппе с образующими $a$ и $b$ и
соотношениями $a^2=\Lambda$, $b^2=\Lambda$, $ab=ba$ (через
$\Lambda$ мы обозначаем пустое слово)? (Ответ: $4$; это группа
$(\bbZ/2\bbZ)\times(\bbZ/2\bbZ)$.)
        \end{problem}

Известно, что существуют такие образующие и соотношения, при
которых проблема равенства слов (выяснить, принадлежат ли два
данных слова одному классу эквивалентности) является
алгоритмически неразрешимой (подробнее
см.~в~\cite{computable-functions}). Мы сейчас покажем, что этот
вопрос можно свести к вопросу о выводимости некоторой формулы в
теории полугрупп, так что если бы она была разрешимой, то
получилось бы противоречие.

Построение такой формулы происходит весьма естественным образом;
мы поясним его на примере. Пусть мы хотим узнать, будут ли
слова~$bb$ и~$a$ равны в полугруппе с образующими~$a$ и~$b$ и
соотношениями $ab\hm=aa$ и~$bab=b$. (Другими словами, мы хотим
узнать, можно ли из слова $bb$ получить слово $a$ с помощью
замен подслов $ab\hm\leftrightarrow aa$ и $bab\hm\leftrightarrow
b$.) Как сформулировать этот вопрос в терминах формул?
Напишем такую формулу:
        $$
\forall a\forall b\,((ab=aa)\land(bab=b)\to(bb=a)).
        $$
Она является теоремой теории полугрупп (истинна во всех
полугруппах, выводима из аксиом полугрупп) тогда и только тогда,
когда слова~$bb$ и~$a$ эквивалентны в указанной полугруппе,
заданной образующими и соотношениями. В самом деле, если одно
слово можно получить из другого заменами, то эти замены
(в~предположении $ab=aa$ и $bab=a$) ничего не меняют и
$bb\hm=a$, так что написанная формула истинна во всех
полугруппах.

Напротив, если слово~$a$ не получается из $bb$ заменой, то
существует
полугруппа, в которой эта формула не истинна: надо взять как раз
полугруппу с образующими~$a$ и~$b$ и соотношениями $ab\hm=aa$ и
$bab\hm=b$, значением переменной~$a$ считать класс слова $a$, а
значением переменной~$b$ считать класс слова~$b$. Тогда
значением терма $ab$ будет класс слова~$ab$, равный классу слова
$aa$ по построению полугруппы. Аналогичным образом при такой
оценке будет истинно и равенство $bab\hm=b$. А равенство $bb=a$
не будет истинно, так как значение терма~$bb$ есть класс
слова~$bb$, значение терма~$a$ есть класс слова~$a$, а эти
классы различны по предположению.

Таким образом, любой алгоритм, проверяющий истинность формул в
классе всех полугрупп, можно было бы использовать для проверки
равенства двух слов в полугруппе, заданной образующими и
соотношениями. А~среди таких полугрупп есть неразрешимые.
        %
\end{proof}

Теория групп (в которой, помимо ассоциативности, есть
ещё аксиомы существования единицы и обратного), также
неразрешима, но доказательство этого сложнее, чем для
полугрупп. Это и не удивительно, поскольку из неразрешимости
теории групп формально выводится неразрешимость теории
полугрупп, как показывает следующая задача.

\begin{problem}
        %
Пусть теория $T$ разрешима, а теория $T'$ той же сигнатуры
получается из~$T$ добавлением конечного числа аксиом. Тогда
теория $T'$ разрешима. (Указание: дополнительные аксиомы
соединяем конъюнкциями и помещаем в посылку импликации.)
        %
\end{problem}

Добавление аксиом может сделать неразрешимую
теорию разрешимой. Например, как мы уже упоминали, это
происходит с теорией групп при добавлении аксиомы
коммутативности.

\subsection*{Теорема Чёрча}%
    \index{Теорема!Чёрча}%
    \glossary{\Чёрч}

Из сказанного легко следует знаменитая теорема Чёрча о
неразрешимости исчисления предикатов:

\begin{theorem}
        \label{church-theorem}
Не существует алгоритма, проверяющего общезначимость формул
первого порядка.
        %
\end{theorem}

В этой формулировке не ограничивается сигнатура (от алгоритма
требуется, чтобы он определял общезначимость формулы с
произвольным числом предикатных и функциональных символов). На
самом деле неразрешимость возникает уже в совсем простых
сигнатурах, как видно из доказательства.

\begin{proof}
        %
Поскольку теория полугрупп конечно аксиоматизируема, то
выводимость формулы $F$ в этой теории равносильна общезначимости
формулы $A\to F$, где $A$\т конъюнкция всех аксиом теории
полугрупп.
        %
\end{proof}

Как мы видим, общезначимость неразрешима уже для формул с
равенством и одним двуместным функциональным символом
(используемым как умножение в полугруппе).

Немного другой вариант рассуждения позволяет доказать
неразрешимость общезначимости для формул с равенством и
одноместными функциональными символами. Заметим, что вопрос о
том, можно ли получить одно слово из другого с помощью замены
подслов по правилам, можно свести к вопросу об общезначимости
некоторой формулы. Пусть, например, мы снова хотим узнать, можно
ли из слова $bb$ получить слово $a$ с помощью замен подслов
$ab\leftrightarrow aa$ и $bab\leftrightarrow b$. Для этого
напишем формулу
        $$
\forall x\,(a(b(x))=a(a(x))) \land
\forall x\,(b(a(b(x)))=b(x))
\to\forall x\,(b(b(x))=a(x)).
        $$
Несложно понять, что её общезначимость равносильна возможности
получить $a$ из $bb$ с помощью указанных замен. В самом деле,
цепочка замен даёт цепочку равенств, соединяющих $b(b(x))$ и
$a(x)$. Напротив, если замены не переводят $bb$ в $a$, то
существует интерпретация, в которой эта формула ложна. А именно,
в качестве носителя этой интерпретации возьмём множество всех
классов слов (в один класс входят слова, получающиеся друг из
друга с помощью замен). На этих классах корректно определены
функции приписывания слева букв $a$ и $b$. Более точно, если
$U$\т один из классов, то через $A(U)$ мы обозначим класс,
содержащий слова $au$ для всех $u\in U$. (Все эти слова лежат в
одном классе: если $u'$ получается из $u$ заменами, то и $au'$
получается из $au$ теми же заменами.) Аналогично мы определяем
функцию $B$ на классах. Легко понять, что $A(B(U))=A(A(U))$ для
любого класса~$U$. В самом деле, если $u$\т слово из $U$, то
слово~$au$ лежит в $A(U)$, слово~$bu$ лежит в $B(U)$, а слова
$abu$ и $aau$ лежат в классах $A(B(U))$ и $A(A(U))$. Но эти
слова переводятся друг в другой заменой $ab\leftrightarrow aa$,
и потому два этих класса совпадают. Аналогичное утверждение
верно и для второй замены. А вот $B(B(U))=A(U)$ верно не всегда:
если в качестве $U$ взять класс пустого слова, то $B(B(U))$
будет классом слова $bb$, а $A(U)$ будет классом слова $a$, и по
предположению классы будут разными. Так что функции $A$ и $B$ на
классах слов показывают, что наша формула не общезначима.

\subsection*{Формальная арифметика}

Рассмотрим множество натуральных чисел с операциями сложения и
умножения и его элементарную теорию
$\Th(\mathbb{N},{=},{+},{\times})$, то есть множество всех
истинных (в натуральном ряду) формул со сложением и умножением.
Это множество, очевидно, полно. Можно доказать
(см.~\cite{computable-functions}), что оно неразрешимо (и, более
того, неарифметично, как говорит теорема Тарского).%
\index{Теорема!Тарского}
\glossary{\Тарский}

Отсюда следует, что теория $\Th(\mathbb{N},{=},{+},{\times})$ не
является конечно аксиоматизируемой. (В самом деле, эта теория
полна и неразрешима, и можно сослаться на
теорему~\ref{complete-decidable}.) Более того, это же
рассуждение показывает, что не существует разрешимого
множества теорем этой теории, из которых бы выводились все
другие теоремы. Отсюда следует, что классическая система аксиом
формальной арифметики, называемая также \emph{арифметикой Пеано}%
\index{Арифметика!Пеано|defin}%
\index{Формальная арифметика|defin}%
\glossary{\Пеано}
(свойства арифметических операций плюс аксиомы индукции), не
может быть полной: существуют истинные формулы, невыводимые в
формальной арифметике. Это утверждение составляет содержание
знаменитой \emph{теоремы Гёделя о неполноте}.%
\index{Теорема!Гёделя о неполноте|defin}%
\glossary{\Гёдель}

\begin{problem}
        %
Покажите, что нельзя добавить к языку теории
$\Th(\mathbb{N},{=},{+},{\times})$ конечное число выразимых
предикатов так, чтобы после этого проходила элиминация
кванторов.%
\index{Элиминация кванторов}
(Указание: арифметическая иерархия не ограничивается
никаким конечным числом уровней.)
        %
\end{problem}

\begin{problem}
        %
Покажите, что элементарная теория целых чисел со сложением и
умножением сводится к элементарной теории натуральных чисел со
сложением и умножением: по замкнутой формуле $\varphi$ со сложением
и умножением можно алгоритмически построить формулу $\varphi'$
с таким свойством: $\varphi$ истинна в~$\mathbb{Z}$ тогда и только
тогда, когда $\varphi'$ истинна в~$\mathbb{N}$. (Указание: целые
числа можно кодировать парами натуральных.)
        %
\end{problem}

\begin{problem}
        %
(Продолжение) Покажите, что верно и обратное: элементарная
теория натуральных чисел сводится к элементарной теории целых
чисел. (Указание. Если бы в целых числах был порядок, это было
бы совсем просто. Чтобы его ввести, можно использовать теорему
Лагранжа о том, что всякое натуральное число представимо в виде
суммы четырёх квадратов.)
        %
\end{problem}

Будет ли теория $\Th(\mathbb{N},{=},{+},{\times})$ категоричной
в счётной мощности? Другими словами, имеет ли она
счётную модель, не изоморфную стандартной? Раньше,
для более простых ситуаций, нам удавалось указать такие модели
явно. Теперь это не удастся, но есть простое общее рассуждение,
устанавливающее существование нестандартной модели. (Оно
аналогично рассуждению, использованному при доказательстве
теоремы~\ref{lowenheim-skolem-2}.)

        \label{nonstandard-n-construction}%
Рассмотрим последовательность формул $E_0(x)$, $E_1(x)$,
$E_2(x),\dots$ с единственным параметром $x$, где $E_i(x)$\т
любая формула, выражающая в стандартной модели свойство~$x=i$.
(Если бы у нас в языке была константа~$1$, можно было бы считать
$E_i(x)$ бескванторной формулой $x=1+1+\ldots+1$, в правой части
которой стоит терм с $i$ единицами.)

Добавим к сигнатуре новую константу $c$ и рассмотрим теорию,
получаемую из $\Th(\mathbb{N},{=},{+},{\times})$ добавлением
счётного семейства формул $\lnot E_i(c)$ (по существу мы
добавляем формулы $c\hm\ne 0,\,c\hm\ne 1,\dots$, записанные
подходящим образом). Любое конечное подмножество полученной
теории имеет модель (возьмём стандартный натуральный ряд и в
качестве~$c$ выберем достаточно большое число). Следовательно
(теорема~\ref{predicate-compactness-theorem} о~компактности), и
        \index{Теорема!компактности}%
вся эта теория совместна. Рассмотрим её счётную нормальную
модель и забудем о символе~$c$; получится некоторая
интерпретация сигнатуры $({=},{+},{\times})$. Она будет
элементарно эквивалентна стандартному натуральному ряду (все
истинные в $\mathbb{N}$ формулы будут истинны по построению, а
все ложные будут ложны, так как их отрицания
истинны).

Осталось показать, что она не будет изоморфной натуральному
ряду. В самом деле, рассмотрим элемент, который является
значением константы~$c$ в нестандартной модели. Он не может
переходить ни в какое натуральное число $i$, поскольку для
соответствующих (друг другу при изоморфизме) элементов выполнены
одни и те же формулы, $E_i(i)$ истинно в натуральном ряду и
$E_i(c)$ ложно в новой интерпретации.

\begin{problem}
        %
Покажите, что найдётся нормальная интерпретация сколь угодно
большой мощности, элементарно эквивалентная натуральным числам
со сложением и умножением.
        %
\end{problem}

\section{Диаграммы и расширения}

В разделе~\ref{lowenheim-skolem-2}
мы видели, что элементарные расширения
интерпретации~$A$ суть модели теории~$\Th_A(A)$. А что можно
сказать о расширениях (без требования элементарности)?
Оказывается, что ситуация тут аналогична, только теория будет
бескванторной.

Пусть дана нормальная интерпретация~$A$ сигнатуры~$\sigma$
(включающей равенство). Как и в прошлом разделе, рассмотрим
сигнатуру~$\sigma_A$, которая получается добавлением к~$\sigma$
констант для всех элементов интерпретации~$A$. Рассмотрим теперь
все бескванторные формулы сигнатуры~$\sigma_A$, истинные в~$A$.
Это множество называется \emph{диаграммой}%
\index{Диаграмма интерпретации|defin}%
\index{Интерпретация!диаграмма|defin}
интерпретации~$A$ и
обозначается $D(A)$.

Всякое расширение $B\supset A$ (в котором $A$ является
подструктурой) является моделью теории $D(A)$. В самом деле,
истинность бескванторных формул из $D(A)$ никак не зависит от
присутствия или отсутствия дополнительных элементов, раз
операции на элементах из~$A$ те же самые). Обратно, любую
модель~$B$ теории~$D(A)$ можно считать расширением
интерпретации~$A$, если отождествить $a\in A$ со значением
соответствующей константы в~$B$. (Как и раньше, различные
элементы $A$ не склеиваются\т формула $a_1\ne a_2$ является
бескванторной.)

Теперь мы готовы дать ответ на такой вопрос. Пусть есть
нормальная интерпретация $A$ сигнатуры~$\sigma$ и некоторая
теория~$T$ (с равенством) этой сигнатуры. В каком случае
существует расширение~$B$ интерпретации~$A$, являющееся
нормальной моделью теории $T$?

\begin{theorem}
        \label{extension-criterion}%
Нормальная интерпретация~$A$ сигнатуры~$\sigma$ может быть
расширена до нормальной модели теории~$T$ (с равенством) тогда и
только тогда, когда все $\Pi_1$\д формулы сигнатуры~$\sigma$,
выводимые из~$T$, истинны в~$A$.%
        %
\end{theorem}

\begin{proof}
        %
Если $\Pi_1$\д формула истинна в некоторой структуре, то она
истинна и в подструктуре (область, по которой пробегают
переменные в кванторах всеобщности, только уменьшается). Если
некоторое расширение~$B$ интерпретации~$A$ является моделью
теории~$T$, то все $\Pi_1$\д формулы, выводимые из~$T$, истинны
в $B$, а потому и в~$A$.

Осталось доказать обратное: если в $A$ истинны все $\Pi_1$\д
следствия формул из $T$, то существует искомое расширение.
Согласно сказанному выше, достаточно доказать, что теория
$D(A) \hm\cup T$ непротиворечива. Если это не так, то из
$T$ выводится некоторая бескванторная формула
$\varphi(a_1,\dots,a_n)$, ложная в $A$. Но в формулы теории $T$
константы $a_1,\dots,a_n$ не входят, поэтому их можно заменить
на свежие переменные $x_1,\dots,x_n$ и вывести формулу
$\varphi(x_1,\dots,x_n)$ и затем $\forall x_1\forall
x_2\ldots\forall x_n\, \varphi(x_1,\dots,x_n)$ Таким образом, мы
нашли $\Pi_1$\д теорему теории~$T$, которая ложна в $A$
(поскольку формула $\varphi(a_1,\dots,a_n)$ ложна), вопреки нашему
предположению.
        %
\end{proof}

Рассмотрим пример из алгебры. Пусть $F$ \т множество с заданной
на нём операцией. В каком случае его можно вложить в
коммутативную группу? Согласно
теореме~\ref{extension-criterion}, для этого необходимо и
достаточно, чтобы в $F$ выполнялись все $\Pi_1$\д следствия
аксиом коммутативной группы (записанных в сигнатуре с
единственной операцией умножения). Некоторые из этих аксиом сами
являются $\Pi_1$\д формулами. Таковы, например, свойства
коммутативности и ассоциативности. Другие аксиомы (существование
единицы и обратного) не лежат в~$\Pi_1$ (например, аксиома о
существовании единицы имеет вид~$\exists e \forall x\dots$).
Поэтому они не обязаны выполняться в~$F$. Но их $\Pi_1$\д
следствия, например, правило сокращения
        $$
\forall x\forall y\forall z\, ({(xy=xz)}\hm\to{(y=z)}),
        $$
должны выполняться. В данном случае оказывается, что этих трёх
утверждений достаточно: всякая коммутативная полугруппа с
сокращением может быть вложена в коммутативную группу.

\begin{problem}
        %
Докажите это утверждение. (Указание. Элементами группы можно
считать классы формальных выражений вида~$x-y$, как это
делается, когда от натуральных чисел переходят к целым. В общей
ситуации эту группу называют группой Гротендика.)
        %
\end{problem}

Вот ещё один хорошо известный пример из алгебры. В каком случае
коммутативное кольцо~$K$ может быть вложено в поле?
Теорема~\ref{extension-criterion} требует, чтобы в~$K$
выполнялись все $\Pi_1$\д теоремы теории полей. Оказывается, что
достаточно выполнения единственного $\Pi_1$\д свойства:
отсутствия делителей нуля:
        $$
\forall x \forall y \, ((xy=0)\to ((x=0)\lor (y=0))).
        $$
В этом случае кольцо может быть вложено в поле.

\begin{problem}
        %
Докажите это утверждение. (Указание. Это поле называют
\emph{полем частных};%
\index{Поле частных|defin}
его элементами являются формальные дроби
вида $m/n$ при естественных определениях равенства и операций.)
        %
\end{problem}

Не всегда, однако, можно указать простые критерии вложимости. Мы
не зря требовали коммутативности: известный советский алгебраист
и логик А.\,И.\,Мальцев%
\glossary{\Мальцев}
доказал, что не всякое некоммутативное
кольцо без делителей нуля вкладывается в тело и что никакое
конечное число $\Pi_1$\д формул не дают критерия вложимости
полугруппы в группу (подробнее см.~в книге Куроша%
\glossary{\Курош}%
~\cite{kurosh},
глава~II, параграф~5).

Мы знаем теперь, когда данную интерпретацию можно расширить до
модели данной теории. Это позволяет легко ответить и на такой
вопрос: когда существует модель данной теории и её расширение,
являющееся моделью другой теории.

\begin{theorem}
        \label{pair-of-theories}%
Пусть даны две теории (с равенством)~$T_1$ и $T_2$ некоторой
сигнатуры. Тогда следующие свойства равносильны:

        (\textbf{а})~%
существует нормальная модель теории $T_1$ и её расширение,
являющееся нормальной моделью теории~$T_2$;

        (\textbf{б})~%
объединение $T_1$ со всеми $\Pi_1$\д теоремами теории $T_2$
совместно;

        (\textbf{в})~%
объединение $T_2$ со всеми $\Sigma_1$\д теоремами теории $T_1$
совместно.
        %
\end{theorem}

\begin{proof}
        %
Прежде всего отметим, что из~(а) очевидно следуют~(б) и~(в). В
самом деле, если $M_1\subset M_2$\т модели соответствующих
теорий, то в $M_1$ истинны все теоремы теории~$T_1$ и все
$\Pi_1$\д теоремы теории~$T_2$ (поскольку они наследуются
из~$M_2$), а в $M_2$ истинны все теоремы теории~$T_2$ и все
$\Sigma_1$\д теоремы теории~$T_1$.

Легко проверить, что симметричные условия~(б) и~(в) равносильны
друг другу, а также такому свойству: не существует $\Sigma_1$\д
теоремы $\exists x_1\ldots\exists x_n\,\varphi$ теории~$T_1$ и
отрицающей её~$\Pi_1$\д теоремы $\forall x_1\ldots\forall
x_n\,\lnot\varphi$ теории~$T_2$. Пусть, например, теория $T_1$
несовместна с $\Pi_1$\д следствиями теории~$T_2$. В этом
противоречии участвует конечное число $\Pi_1$\д формул,
которые можно объединить в одну
(см.~раздел~\ref{prenex-normal-form},
с.~\pageref{sigma-pi-closed}). Получится $\Pi_1$\д формула, она
будет выводима в теории~$T_2$, а её отрицание выводимо в $T_1$.

Нам осталось доказать, что любое из свойств~(б) и~(в) влечёт~(а).
Здесь нам придётся нарушить симметрию и использовать именно~(б).
По условию есть интерпретация $M_1$, в которой истинны все
теоремы теории~$T_1$ и все $\Pi_1$\д теоремы теории~$T_2$.
Согласно теореме~\ref{extension-criterion} найдётся её
расширение~$M_2$, являющееся моделью~$T_2$, что и требовалось
доказать.
        %
\end{proof}

Можно было бы пытаться рассуждать симметричным образом,
начав с модели теории~$T_2$, в которой истинны все $\Pi_1$\д теоремы
теории~$T_1$, и пытаться выделить в ней подструктуру, являющуюся
моделью теории~$T_1$. Однако этот план не проходит, поскольку
аналог теоремы~\ref{extension-criterion} для подструктур неверен.

\begin{problem}
        %
Покажите, что возможна такая ситуация: все $\Sigma_1$\д теоремы
некоторой теории~$T$ истинны в некоторой интерпретации~$M$, но
$M$ не имеет подструктуры, являющейся моделью теории~$T$.
(Указание. Рассмотрим теорию линейно упорядоченных множеств без
минимального элемента. Все её $\Sigma_1$\д следствия верны в
$\mathbb{N}\hm+\mathbb{Z}$, поскольку переносятся
из~$\mathbb{Z}$, поэтому в силу элементарной эквивалентности
верны и в~$\mathbb{N}$.)
        %
\end{problem}

Вот ещё одно следствие доказанных в этом разделе результатов.
Теорию~$T$ называют~\emph{$\Pi_1$\д аксиоматизируемой},%
\index{$\Pi_1$\д аксиоматизируемая теория|defin}%
\index{Теория!$\Pi_1$\д аксиоматизируемая|defin}
если существует множество $\Pi_1$\д формул, из которого
выводятся все теоремы теории~$T$ и только они.

Напомним, что нормальная интерпретация $A$ сигнатуры $\sigma$
        \index{Подструктура}%
является \emph{подструктурой} нормальной интерпретации $B$ той
же сигнатуры, если $B$ является расширением $A$, то есть
носитель интерпретации $A$ есть подмножество носителя
интерпретации~$B$ и функциональные и предикатные символы
интерпретируются одинаково на аргументах из $A$.
(Другими словами, чтобы задать какую-либо подструктуру данной
нормальной интерпретации $B$, нужно выбрать подмножество
носителя $B$, замкнутое относительно сигнатурных операций.)

\begin{theorem}[Лося\ч Тарского]
\index{Теорема!Лося\ч Тарского|defin}%
\glossary{\Лось}%
\glossary{\Тарский}%
        %
Теория $\Pi_1$\д аксиоматизируема тогда и только тогда, когда
она устойчива относительна перехода к подструктурам, то есть
        \index{Устойчивость!при переходе к подструктурам}%
когда любая подструктура любой её нормальной модели является её
моделью.
        %
\end{theorem}

\begin{proof}
        %
Очевидно, $\Pi_1$\д аксиоматизируемая теория устойчива
относительно перехода к подструктурам (все формулы из её
$\Pi_1$\д аксиоматизации остаются истинными). Об\-рат\-но, пусть
$T$\т произвольная теория, устойчивая относительно перехода к
подструктурам. Рассмотрим множество $T_1$ всех $\Pi_1$\д формул,
выводимых в~$T$. Проверим, что все теоремы $T$ выводятся
из~$T_1$. Пусть какая\д то формула $\varphi$ выводится из~$T$,
но не из~$T_1$. Тогда теория~$T_1+\lnot\varphi$ непротиворечива
и по теореме~\ref{pair-of-theories} найдётся (нормальная) модель
теории~$\{\lnot\varphi\}$ и её расширение, являющееся моделью
теории~$T$, что противоречит предположению.
        %
\end{proof}

\begin{problem}
        %
Докажите, что если формула устойчива относительно перехода к
подструктурам, то она выводимо эквивалентна некоторой $\Pi_1$\д формуле той
же сигнатуры.
        %
\end{problem}

Симметричное рассуждение доказывает симметричное утверждение про
$\Sigma_1$\д аксиоматизируемые теории.%
\index{$\Sigma_1$\д аксиоматизируемая теория|defin}%
\index{Теория!$\Sigma_1$\д аксиоматизируемая|defin}

\begin{theorem}
        %
Теория является~$\Sigma_1$\д аксиоматизируемой тогда и только
        \index{Устойчивость!при расширении}%
тогда, когда она устойчива относительно перехода к рас\-ши\-ре\-ниям.
        %
\end{theorem}

\begin{problem}
        %
Проведите подробно соответствующее рассуждение (дав необходимые
определения).
        %
\end{problem}

\begin{problem}
        %
Докажите, что если формула устойчива относительно перехода к
расширениям, то она выводимо эквивалентна некоторой $\Sigma_1$\д формуле
той же сигнатуры.
        %
\end{problem}

Теоретико\д модельные критерии существуют и для других классов
формул, в частности $\Pi_2$\д формул (то есть формул типа
$\forall\exists$). Такие формулы не устойчивы ни относительно
расширений, ни относительно подструктур. Рассмотрим, например,
утверждение об отсутствии наибольшего элемента в упорядоченном
множестве. Оно записывается в виде $\forall\exists$\д формулы.
Истинность его в некотором множестве вовсе не влечёт его
истинность в подмножествах и в расширениях. Тем не менее кое\д
что об этом утверждении сказать можно: если ни одно из множеств
возрастающей цепи
        $
M_0\hm\subset M_1\hm\subset M_2\hm\subset\dots
        $
не имеет наибольшего элемента, то и объединение~$\cup_i M_i$ не
имеет наибольшего элемента (проверьте). Именно это свойство, как
мы вскоре увидим, характеризует $\Pi_2$\д формулы.

Пусть дана последовательность
        $$
M_0\hm\subset M_1\hm\subset M_2\hm\subset\ldots
        $$
нормальных (в этом разделе мы другие не рассматриваем)
интерпретаций сигнатуры~$\sigma$, причём~$M_i$ является
подструктурой~$M_{i+1}$ (предикаты и функции согласованы). Тогда
объединение этой возрастающей цепи интерпретаций также является
(нормальной) интерпретацией сигнатуры~$\sigma$. (Подобная
конструкция используется в теории полей, когда строится
алгебраическое замыкание счётного поля: мы расширяем поле,
добавляя по очереди корни различных многочленов, а потом берём
объединение этих полей.)

Заметим, что любая $\Pi_2$\д формула устойчива относительно
объединения цепей: если она истинна во всех $M_i$, то она
истинна и в их объединении. В самом деле, пусть формула $\forall
x \exists y\, \varphi(x,y)$ с бескванторной частью
$\varphi(x,y)$ истинна во всех~$M_i$. Тогда она истинна и в их
объединении. В самом деле, любое $x$ из объединения принадлежит
какому\д то~$M_i$, и в том же самом~$M_i$ можно найти
подходящее~$y$. (Если переменных несколько, рассуждение
аналогично.)

Поэтому и любая теория, имеющая $\Pi_2$\д аксиоматизацию,
        \index{Устойчивость!относительно объединения цепей}%
устойчива относительно объединения. Обратное утверждение также
верно:

\begin{theorem}[Чэна\ч Лося\ч Сушко]
\index{Теорема!Чэна\ч Лося\ч Сушко|defin}%
\glossary{\Чэн}%
\glossary{\Лось}%
\glossary{\Сушко}%
        %
Теория является $\Pi_2$\Д аксиоматизируемой тогда и только
тогда, когда она устойчива относительно объединения возрастающих
цепей (объединение любой цепи её моделей также является
её моделью).
        %
\end{theorem}

\begin{proof}
        %
Доказательство этой теоремы использует понятие элементарного
расширения. Напомним, что $M_2$ называется элементарным
расширением $M_1$, если $M_1\hm\subset M_2$ и в $M_2$ истинны
те же формулы с константами из $M_1$, что и в~$M_1$.
(Обозначение: $M_1\hm\prec M_2$.)

\begin{problem}
        %
Покажите, что если $M_1\hm\prec M_2 \hm\prec M_3$, то $M_3$
есть элементарное расширение~$M_1$.
        %
\end{problem}

\textsf{Лемма Тарского}.%
\index{Лемма!Тарского|defin}%
\glossary{\Тарский}%
\index{Цепь элементарных расширений}
Объединение цепи элементарных расширений $M_1\hm\prec
M_2\hm\prec M_3\hm\prec\ldots$ является элементарным расширением
каждой из интерпретаций цепи.

Доказательство леммы. Пусть параметрам формулы~$\varphi$ приданы
значения в каком\д либо из~$M_i$. Нам надо доказать, что
полученная формула одновременно истинна или ложна в $M_i$ и в
объединении цепи, которое мы обозначим через~$M$. (Условие
леммы гарантирует, что формула~$\varphi$ с указанными значениями параметров
одновременно истинна или ложна во всех интерпретациях цепи, начиная с $M_i$.)

Это утверждение доказывается индукцией по построению
формулы~$\varphi$. Для атомарных формул оно очевидно; для
логических операций индукция также проходит автоматически.
Единственный содержательный случай\т это кванторы. Пусть формула
$\varphi$ начинается с квантора~$\exists \xi$. Если подходящее
значение $\xi$ найдётся уже в $M_i$, то оно годится и для~$M$
(пользуемся предположением индукции). В обратную сторону: если
подходящее $\xi$ найдётся в~$M$, то оно принадлежит $M_j$ при
достаточно большом~$j$, поэтому формула истинна в $M_j$
(предположение индукции). Остаётся вспомнить, что~$M_j$
элементарно эквивалентно~$M_i$.

Как всегда, квантор всеобщности можно выразить с помощью квантора
сушествования (или провести двойственное рассуждение).
Лемма
Тарского доказана.

Теперь докажем теорему Чэна\Ч Лося\Ч Сушко. Предположим, что
теория~$T$ устойчива относительно объединения цепей. Обозначим
через $T'$ множество всех $\Pi_2$\д теорем $T$. Нам надо
доказать, что любая модель $T'$ является  моделью $T$.

Для этого, начав с любой модели $M'$ теории $T'$, мы построим
цепь интерпретаций
         $$
M'= M_0\subset M_1\subset M_2\subset M_3\subset\ldots,
         $$
в которой чередуются модели теории~$T'$ (стоящие на чётных местах интерпретации
$M_0,M_2,M_4,\dots$), которые являются элементарными
расширениями друг друга, и модели теории $T$ (интерпретации
$M_1,M_3,M_5,\dots$; они, впрочем, также будут моделями
теории~$T'$).

Объединение всех~$M_i$ будет моделью теории~$T$, так как эта
теория устойчива относительно расширений. С другой стороны, по
лемме Тарского это объединение элементарно эквивалентно
интерпретациям $M_0,M_2,M_4,\dots$ Поэтому все они, включая
исходную модель $M'\hm=M_0$, будут моделями теории~$T$, что и
требовалось доказать.

Осталось построить требуемую цепь. Интерпретация $M_0\hm=M'$ уже
есть. Будем строить цепь по шагам, продолжая её на каждом шаге
на два звена вперёд. Возможность этого обеспечивает такая лемма:

\textsf{Лемма о расширении}.%
\index{Лемма!о расширении}
Если все $\Pi_2$\д следствия теории~$T$ истинны в интерпретации
$A$, то можно построить её расширения $A\subset B \subset C$ так,
чтобы $B$ было моделью теории~$T$, а $C$ было элементарным
расширением~$A$.

Прежде чем доказывать лемму о расширении, покажем (хотя это нам
и не понадобится), что сформулированное условие необходимо.
Пусть $A\hm\subset B\hm\subset C$, причём $C$\т элементарное
расширение~$A$. Тогда любое $\Pi_2$\д утверждение, истинное
в~$B$, истинно и в~$A$. В самом деле, пусть утверждение $\forall
x \exists y \varphi(x,y)$ с бескванторной формулой~$\varphi$
истинно в $B$. Проверим его истинность в $A$. Если оно ложно при
$x=a$, то $\exists y\,\varphi(a,y)$ ложно и в~$A$, и в $C$
(элементарность расширения) и потому не может быть истинным
в~$B$ (поскольку всякое~$y$ из~$B$ лежит и в~$C$).

Доказательство леммы о расширении. Что требуется от данного
расширения~$B$ интерпретации~$A$, чтобы можно было построить~$C$
с требуемыми свойствами? Свойства эти состоят в том, что
$C$~должно быть моделью теории $\Th_A(A)$ и расширением
интерпретации~$B$. Как раз про это говорит
теорема~\ref{extension-criterion}, надо лишь в качестве $\sigma$
в этой теореме взять сигнатуру $\sigma$ с добавленными
константами для $A$ (мы обозначали её $\sigma_A$), а в качестве
теории $T$ из теоремы~\ref{extension-criterion} взять
$\Th_A(A)$, то есть множество всех истинных в~$A$ формул с
константами из~$A$.

Применяя указанный в теореме~\ref{extension-criterion} критерий,
можно сформулировать утверждение, которое нам осталось доказать,
так: найдётся модель~$B$ теории~$T$, которая является
расширением~$A$ и в которой истинны все $\Pi_1$\д формулы
сигнатуры $\sigma_A$, выводимые из $\Th_A(A)$. Вспоминая метод
диаграмм, можно сказать, что нас интересует совместность
теории~$T$ с $D(A)$ и со всеми $\Pi_1$\д следствиями теории
$\Th_A(A)$ в сигнатуре $\sigma_A$. В данном случае $D(A)$ можно и
не упоминать явно, так как оно содержится в $\Th_A(A)$.

Итак, докажем, что теория $T$ совместна со всеми
$\Pi_1$\д формулами с константами из~$A$, истинными в~$A$. Если
это не так, из~$T$ выводится отрицание какой\д то из этих
формул, то есть некоторая $\Sigma_1$\д формула
        $$
\exists \beta_1\dots\exists \beta_m\,
           \lnot\varphi(a_1,\dots,a_n,\beta_1,\dots,\beta_m),
        $$
ложная в $A$. Константы $a_1,\dots,a_n$ не входят в теорию~$T$,
поэтому из $T$ выводится и формула
        $$
\forall \alpha_1\ldots\forall \alpha_n
\exists \beta_1\dots\exists \beta_m\,
           \lnot\varphi(\alpha_1,\dots,\alpha_n,\beta_1,\dots,\beta_m),
        $$
которая будет выводимой из~$T$ формулой класса~$\Pi_2$, ложной в~$A$,
а таких формул не бывает по условию.

Лемма о расширении (а с ней и теорема Чэна\ч Лося\ч Сушко) доказана.
        %
\end{proof}

\begin{problem}
        %
Докажите, что если формула устойчива относительно объединения
возрастающих цепей, то она выводимо эквивалентна
$\Pi_2$\д формуле той же сигнатуры.
        %
\end{problem}

\begin{problem}
        %
Покажите, что две интерпретации одной сигнатуры элементарно
эквивалентны тогда и только тогда, когда они имеют изоморфные
элементарные расширения.
        %
\end{problem}



\section{Ультрафильтры и компактность}
        \label{ultraproducts}

В этом разделе мы дадим прямое доказательство теоремы о
        \index{Теорема!компактности, прямое доказательство}%
компактности (теорема~\ref{predicate-compactness-theorem},
с.~\pageref{predicate-compactness-theorem}).

Пусть $S$\т непустое множество,
а $F$\т непустое семейство подмножеств множества~$S$.
Семейство $F$ называется \emph{фильтром}%
\index{Фильтр|defin}
на $S$, если выполнены следующие
три свойства (для наглядности
множества из $F$ мы называем далее \emph{большими}):%
\index{Большое множество|defin}%
\index{Множество!большое|defin}
\begin{itemize}
        %
\item $\varnothing\not\in F$
(пустое множество не является большим);
        %
\item $A\in F,\ B\in F\Rightarrow A\cap B\in F$
(пересечение двух больших множеств\т снова большое множество);
отсюда следует, что пересечение конечного числа больших
множеств\т большое множество и что любые два больших
множества пересекаются;
        %
\item $A\in F,\ A\subset B\Rightarrow B\in F$ (любое
надмножество большого множества является большим); отсюда
следует, что всё множество $S$ является большим.
        %
\end{itemize}

Дополнения к большим множествам естественно назвать
\emph{малыми}.%
\index{Малое множество|defin}%
\index{Множество!малое|defin}
(Отметим, что множество может не быть ни большим,
ни малым с точки зрения фильтра; это отличает фильтры от ультрафильтров, которые мы
вскоре определим.)

Тривиальным примером фильтра является семейство, состоящее из
единственного множества $S$. Другой пример\т семейство всех
подмножеств $S$, содержащих некоторый выделенный элемент
$s\hm\in S$. Такой фильтр называется \emph{главным}.%
\index{Главный фильтр|defin}%
\index{Фильтр!главный|defin}
Третий
пример: пусть $S$ бесконечно, тогда фильтром будет множество
всех \emph{коконечных}%
\index{Коконечное подмножество|defin}%
\index{Подмножество!коконечное|defin}
подмножеств $S$, то есть подмножеств,
дополнение которых конечно. (Другими словами, малыми будут
конечные множества.) Последний пример\т из анализа: фильтром
на~$\mathbb{R}$ является семейство всех окрестностей некоторой
фиксированной точки~$a$, то есть всех множеств, для которых~$a$
является внутренней точкой.

\begin{problem}
        %
Переформулируйте определение фильтра в терминах свойств малых множеств.
        %
\end{problem}

Заметим, что одно и то же множество не может быть одновременно и
большим, и малым, поскольку любые два больших множества
пересекаются (по большому, и потому непустому, множеству). Но,
как мы уже говорили, множество может не быть ни большим, ни
малым. Если таких \лк промежуточных\пк\ множеств нет, фильтр
называют ультрафильтром.

Иными словами, \emph{ультрафильтром}%
\index{Ультрафильтр}
называется фильтр на~$S$ с
таким свойством: $A\hm\in F$ или ${S\setminus A}\hm\in F$ для
любого множества~$A\hm\subset S$.

Очевидно, любой главный фильтр является ультрафильтром. В
остальных наших примерах фильтры не были ультрафильтрами.

\begin{problem}
        \label{finite-ultrafilter}%
Докажите, что на конечном множестве любой ультрафильтр является
главным.
        %
\end{problem}

\begin{problem}
        %
Докажите, что любой неглавный ультрафильтр содержит все
коконечные множества.
        %
\end{problem}

\begin{problem}
        \label{finite-invariant}%
Докажите, что если ультрафильтр не является главным, то вместе
с каждым множеством~$A$ он содержит и все множества~$B$, для
которых симметрическая разность $A\bigtriangleup B$ конечна.
        %
\end{problem}

Определение ультрафильтра можно переформулировать следующим
образом: ультрафильтры\т это фильтры, не имеющие собственных
расширений (максимальные по включению). Докажем это. Если фильтр
$F$ не максимален, то найдётся больший фильтр $F'$. Тогда
множество $A\hm\in F'\setminus F$ и его дополнение до $S$ не
принадлежит~$F$ (иначе и $A$, и его дополнение принадлежали бы
фильтру~$F'$). Следовательно, $F$ не является ульт\-ра\-фильтром.

Обратно, пусть фильтр $F$ не является ультрафильтром, и ни
множество $A$, ни его дополнение~$S\setminus A$ не принадлежат
$F$. Добавим к~$F$ все множества вида $A\cap B$ (для
всех~$B\hm\in F$) и все их надмножества. Получится фильтр. В
самом деле, пустое множество ему не принадлежит, так как иначе
бы $A$ не пересекалось с некоторым множеством из~$F$ и
$S\hm\setminus A$ содержало бы некоторое множество из~$F$ потому
лежало бы в~$F$. Остальные свойства фильтра очевидны; новый
фильтр (в отличие от исходного) содержит~$A$ и потому
расширяет~$F$.

\begin{theorem}
        \label{filter-ultrafilter}%
Всякий фильтр~$F$  на множестве $S$ можно расширить до ультрафильтра
$F'\hm\supset F$.
        %
\end{theorem}

\begin{proof}
        %
Доказательство этой теоремы неконструктивно: мы не предъявляем
такого фильтра, а устанавливаем его существование с помощью
леммы Цорна%
\index{Лемма!Цорна}%
\glossary{\Цорн}
(см.~\cite{naive-set-theory}).
Нужно только заметить, что объединение любой цепи фильтров
является фильтром (что непосредственно следует из определения).

Другими словами, пока фильтр не станет ультрафильтром, мы берём
\лк промежуточное\пк\ множество и расширяем фильтр, объявляя его
большим, повторяя этот процесс по трансфинитной индукции.
        %
\end{proof}

\begin{problem}
        %
Докажите, что на любом бесконечном множестве есть неглавный
ультрафильтр. (Указание: расширим фильтр коконечных множеств до
ультрафильтра.)
        %
\end{problem}

Можно представлять себе элементы множества~$S$ как голосующих
(которые никогда не воздерживаются от голосования). При этом
фильтр на~$S$ определяет регламент: решение принимается, если
множество проголосовавших \лк за\пк\т большое. Аксиомы
фильтра тогда звучат так: решение, против которого все, принято
быть не может; если каждое из двух решений принимается, то они
принимаются и в совокупности; наконец, принятое решение не может
быть отвергнуто, если некоторые из голосовавших против него
передумали.

Свойство ультрафильтра также имеет ясный смысл: по любому
вопросу можно принять решение (одно из двух противоположных
мнений набирает большинство). Главные ультрафильтры
соответствуют диктатуре (существенно мнение лишь одного
голосующего); задача~\ref{finite-ultrafilter} показывает, что
для конечного числа голосующих любые другие способы не позволяют
принять решения по некоторым вопросам.

\smallskip

Ультрафильтры можно использовать для построения любопытных
примеров. Вот один из них. Рассмотрим игру двух участников, в
которой они по очереди объявляют некоторые натуральные числа \лк
своими\пк. На первом шаге начинающий игру объявляет своими числа
от нуля до некоторого числа $n_1$, на втором шаге его противник
присваивает числа от $n_1$ (не включая его) до некоторого
большего числа $n_2$ (включая его), затем первый игрок
присваивает числа от $n_2$ до $n_3$ и так далее. Партия
продолжается неограниченно и делит натуральный ряд между первым
и вторым (на два взаимно дополнительных множества). Выигрывает
тот, чьё множество большое (принадлежит некоторому фильтру).

Если этот фильтр является ультрафильтром, то в этой игре не
может быть ничьей. Если ультрафильтр главный, то игра
тривиальна\т побеждает тот, кто захватит решающее число, и
потому первый может гарантировать выигрыш на первом же ходу.

\begin{theorem}
        %
Если ультрафильтр неглавный, то ни один из игроков не имеет
выигрышной стратегии. (\emph{Стратегия}%
\index{Стратегия|defin}%
\т это функция,
предписывающая следующий ход в зависимости от истории игры.
Стратегия считается \emph{выигрышной},%
\index{Выигрышная стратегия|defin}%
\index{Стратегия!выигрышная|defin}
если её использование
гарантирует выигрыш при любой игре противника.)
        %
\end{theorem}

\begin{proof}
        %
Прежде всего отметим, что оба игрока не могут одновременно
иметь выигрышные стратегии. (Что будет, если они оба ими
воспользуются?) Покажем теперь, что если выигрышную стратегию
имеет один, то её имеет и второй. Совсем просто понять, что если
у ходящего вторым есть выигрышная стратегия, то и первый может
ей воспользоваться (он должен представить себя вторым, считая,
что первый на первом ходу ничего не взял).

Не столь ясно, что выигрышная стратегия первого может быть
использована вторым, но и это верно\т поскольку конечные
множества не влияют на принадлежность ультрафильтру
(задача~\ref{finite-invariant}), второй может забыть про ход, с
которого началась игра, и вообразить себя первым. (Это сделает
его первый ход бессмысленным, если этот ход окажется меньше хода
противника. В этом случае можно сделать сразу второй ход первого игрока,
и далее следовать стратегии.)
        %
\end{proof}

\begin{problem}
        %
Проведите это рассуждение подробно.
        %
\end{problem}

\smallskip

Сейчас мы докажем теорему компактности с помощью ультрафильтров.
Для этого нам понадобится понятие произведения интерпретаций.

Пусть $M_s$\т семейство интерпретаций некоторой (одной и той же)
сигнатуры~$\sigma$, индексированное множеством~$S$ (для каждого
$s\hm\in S$ имеется своя интерпретация~$M_s$). Определим
произведение интерпретаций этого семейства, обозначаемое
        $$
\prod_{s\in S} M_s.
        $$
Элементами носителя будут отображения, сопоставляющие c каждым
индексом $s\hm\in S$ некоторый элемент интерпретации~$M_s$.
Иными словами, носитель строимой интерпретации будет декартовым
произведением всех $M_s$.

Функциональные символы интерпретировать легко: они применяются
отдельно в каждой компоненте. Именно так определяется
произведение групп или колец в алгебре. Остаётся определить
предикатные символы. В алгебре два элемента в произведении колец
или групп считаются равными, если все их компоненты равны. По
аналогии будем считать, что два элемента $s\mapsto a_s$ и
$s\mapsto b_s$ делают истинным двуместный предикат $P$, если
$P(a_s,b_s)$ истинно в интерпретации $M_s$ для всех $s$. (Мы
взяли двуместный символ для примера, то же самое можно сделать и
для символов любой валентности.)

Таким способом в произведении двух упорядоченных множеств (индексное
множество~$S$ равно $\{1,2\}$, сигнатура есть $({=},{\le})$)
возникает покомпонентный порядок на парах: $\langle
a_1,a_2\rangle\hm\le\langle b_1,b_2\rangle$, если $a_1\hm\le
a_2$ и $b_1\hm\le b_2$. Заметим, что такой порядок на
произведении двух линейно упорядоченных множеств уже не будет
линейным (если сомножители состоят более чем из одного
элемента).

Нам это не нравится: мы хотим, чтобы произведение интерпретаций
обладало бы всеми свойствами, которыми обладают сомножители.
Введём понятие \emph{фильтрованного произведения}%
\index{Фильтрованное произведение|defin}
(по модулю
данного фильтра). Пусть на множестве индексов $S$ задан фильтр
$F$. Изменим определение истинности предикатов и будем считать,
что элементы $s\hm\mapsto a_s,\,s\hm\mapsto b_s,\dots$ делают
истинным предикат~$P$, если $P(a_s,b_s,\dots)$ истинно \лк для
большинства~$s$\пк, то есть если множество $\{s\mid
P(a_s,b_s,\dots)\}$ принадлежит фильтру~$F$. В остальном
(носитель, функциональные сим\-волы) определение остаётся прежним.

Что будет равенством в фильтрованном произведении нормальных
интерпретаций? Два элемента произведения (то есть функции на
множестве индексов) равны, если они \лк совпадают почти
всюду\пк, то есть множество индексов, где они совпадают,
принадлежит фильтру. При этом полученная интерпретация не будет
нормальной. Чтобы перейти к нормальной, нужно рассмотреть классы
равных элементов\т как это делается, скажем, для пространства
$L_2$ интегрируемых с квадратом функций, где элементами являются
не сами функции, а их классы с точностью до совпадения почти
всюду.

\begin{problem}
        %
Что можно сказать про фильтрованное произведение по главному
фильтру?
        %
\end{problem}

Вернёмся к нашему примеру: произведению линейно упорядоченных
множеств. Будет ли оно линейно упорядоченным? Это зависит от
фильтра. Например, если фильтр состоит только из множества $S$,
то фильтрованное произведение совпадает с определённым ранее, и
линейного порядка не получится. Но если фильтр является
ультрафильтром, то будет. В самом деле, рассмотрим два элемента
$s\hm\mapsto a_s$ и $s\hm\mapsto b_s$ в произведении и два
множества $\{s\mid a_s\hm\le b_s\}$ и $\{s\mid a_s\hm\ge b_s\}$.
В объединении они покрывают всё~$S$, и потому (если у нас
ультрафильтр) одно из них должно быть большим (если оно не
большое, то оно малое, его дополнение большое и содержится во
втором множестве).

\begin{problem}
        %
Докажите, что в фильтрованном произведении нормальных
интерпретаций функции и предикаты корректны относительно
равенства (то есть совпадения почти всюду): при замене
аргументов на равные значение функции совпадает с прежним почти
всюду, а значение предиката не меняется.
        %
\end{problem}

Это утверждение можно сформулировать так: аксиомы равенства
истинны в фильтрованном произведении нормальных интерпретаций.
Для ультрафильтров верно и более общее свойство: любая формула,
истинная во всех интерпретациях, истинна в фильтрованном
произведении по модулю ультрафильтра. Поэтому именно такие
\emph{ультрапроизведения}%
\index{Ультрапроизведение|defin}
(фильтрованные по модулю
ультрафильтра) представляют основной интерес для логики.

Мы сейчас докажем это свойство по индукции. Как обычно, надо
предварительно распространить его на формулы с параметрами.

\begin{theorem}[Лося об ультрапроизведениях]
\index{Теорема!Лося об ультрапроизведениях}%
\glossary{\Лось}%
        %
Пусть параметрами формулы $\varphi$ являются переменные
$a,b,\dots$ Она будет истинной в ультрапроизведении $\prod_{s\in
S} M_s$ при значениях параметров $\alpha,\beta,\dots$ тогда и
только тогда, когда множество тех~$s$, при которых $\varphi$
истинна в $M_s$ при значениях параметров
$\alpha_s,\beta_s,\dots$, принадлежит ульт\-ра\-фильтру.
        %
\end{theorem}

Наглядно утверждение теоремы можно сформулировать так:
голосование можно проводить не только по атомарным вопросам, а
для любых формул. Для замкнутых формул про параметры можно
ничего не говорить, и мы получаем, что формула истинна в
ультрапроизведении, если и только если она истинна в большинстве
(с точки зрения ультрафильтра) сомножителей. В частности, если формула
истинна во всех сомножителях, то она истинна и в ультрапроизведении.
Это важное утверждение заслуживает особого упоминания:

\textsf{Следствие}. Ультрапроизведение семейства
моделей некоторой те\-о\-рии
является моделью той же теории.

\begin{proof}
Докажем  теорему Лося индукцией по построению формулы.
Для атомарных формул оно непосредственно следует из определения
истинности предикатов.

Пусть формула $\varphi$ является конъюнкцией двух других формул
$\psi\land\eta$, для которых утверждение уже верно. Тогда
множество тех индексов, для которых $\varphi$ истинно, является
пересечением множеств тех индексов, где истинны~$\psi$ и~$\eta$.
Тем самым нам нужно такое свойство ультрафильтра: пересечение
двух множеств является большим тогда и только тогда, когда оба
они большие. Оно непосредственно следует из определения фильтра
(здесь неважно, что это ультрафильтр), поскольку пересечение
содержится в обоих множествах.

Для объединения соответствующее свойство звучит так: объединение
$S\cup T$ двух множеств большое тогда и только тогда, когда хотя
бы одно из множеств~$S$ и~$T$ большое. В одну сторону (если одно
из множеств большое, то и объединение таково) это вытекает из
определения фильтра. В обратную сторону надо воспользоваться
свойствами ультрафильтра: если $S\hm\cup T$ большое, а оба
множества~$S$ и~$T$\т нет, то они малые, их дополнения большие,
пересечение дополнений большое, и не пересекается с $S\hm\cup
T$, что невозможно.

Пусть формула $\varphi$ имеет вид $\lnot\psi$. Тогда имеет место
такая цепочка: ($\varphi$~истинна в
ультрапроизведении$)\hm\Leftrightarrow(\psi$~ложна в
нём$)\hm\Leftrightarrow($множество индексов тех сомножителей,
где $\psi$ истинна, не является большим$)\hm\Leftrightarrow($это
множество является малым$)\hm\Leftrightarrow($его дополнение
большое$)\Leftrightarrow($мно\-жест\-во номеров тех сомножителей,
где~$\psi$~ложна (то есть~$\varphi$~истинна), большое).

Импликация сводится к уже рассмотренным случаям ($\psi\hm\to\eta$
эквивалентна $\lnot\psi\hm\lor\eta$); можно также сразу заменить
формулу на эквивалентную без импликации.

Наиболее интересен случай кванторов. При этом можно ограничиться
квантором существования (квантор всеобщности сводится к нему и к
отрицаниям). Он разбирается так (мы используем не вполне
корректные обозначения\т надеемся, они не вызовут путаницы).
Пусть формула
        $
\exists x\,\varphi(x,y,z,\dots)
        $
истинна в ультрапроизведении. Это значит, что существует функция
$x\colon s\hm\mapsto x_s$, для которой $\varphi(x,y,z,\dots)$
истинно в ультрапроизведении. По предположению индукции это означает,
что для большинства~$s$ формула $\varphi(x_s,y_s,z_s,\dots)$
истинна в $M_s$. Но тогда для этих индексов~$s$ и формула
        $
\exists x\,\varphi(x,y,z,\dots)
        $
истинна в $M_s$, что и требовалось. Обратное рассуждение
аналогично: если для большинства~$s$ найдётся
соответствующее значение~$x_s$, то эти~$x_s$ можно собрать в функцию
(доопределив её как угодно на малом множестве остальных~$s$), и
эта функция будет искомым значением $x$ в
ультрапроизведении.

Теорема Лося доказана.
        %
\end{proof}

Мы уже говорили, что произведение нормальных интерпретаций может
не быть нормальным. Но теорема Лося гарантирует, что в
ультрапроизведении нормальных интерпретаций выполнены аксиомы
равенства (поскольку они выполнены в каждом сомножителе), и
потому равенство является отношением эквивалентности, и классы
эквивалентности уже дают нормальную интерпретацию (с теми же
истинными формулами), как мы видели в
разделе~\ref{equality-axioms}.

Теорема Лося позволяет дать прямое доказательство теоремы
компактности (теорема~\ref{predicate-compactness-theorem},
        \index{Теорема!компактности, прямое доказательство}%
с.~\pageref{predicate-compactness-theorem}). Она утверждает, что
если всякое конечное подмножество данного множества замкнутых
формул $T$ совместно (имеет модель), то и всё множество $T$
совместно.

Модель для всего множества~$T$ строится как ультрапроизведение.
Индексами будут конечные подмножества множества~$T$. Для каждого
из них сомножителем будет существующая по условию модель. Теперь
надо правильно подобрать фильтр на семействе конечных
подмножеств множества~$T$. Нам нужно, чтобы для каждого~$t\in
T$ семейство всех конечных подмножеств, содержащих~$t$, было бы
большим. (В этом случае теорема Лося гарантирует, что $t$ будет
истинно в ульт\-ра\-про\-из\-ведении.)

Как построить такой фильтр? Для каждого конечного $T'\hm\subset
T$ рассмотрим семейство $S(T')$ всех конечных подмножеств,
содержащих $T'$. Очевидно, пересечение таких семейств снова
будет семейством такого вида ($S(T')\hm\cap S(T'')=S(T'\hm\cup
T'')$), так что после добавления всех надмножеств всех таких
множеств получится фильтр. Остаётся расширить этот фильтр до
ультрафильтра по теореме~\ref{filter-ultrafilter}.
Теорема компактности доказана.

Поучительно проследить до конца, что даёт такого рода построение
для какого\д нибудь конкретного примера. Вспомним построение
нестандартного натурального ряда на
с.~\pageref{nonstandard-n-construction}. Оно использовало теорему
компактности. Сочетая его с приведённым только что доказательством
теоремы компактности (и кое\д что упростив), получаем такую
конструкцию.

Рассмотрим натуральные числа как интерпретацию сигнатуры
$({=},{+},{\times})$. Рассмотрим ультрапроизведение $\nstbb{\bbN}$
счётного числа таких интерпретаций по модулю какого\д либо
неглавного ультрафильтра. Теорема Лося говорит, что в этой
интерпретации будут истинны те же формулы, что в натуральном
ряду, то есть что $\nstbb{\bbN}$ элементарно эквивалентна стандартной
интерпретации $\bbN$.

Покажем, что~$\bbN$ не изоморфна~$\nstbb{\bbN}$. В самом деле, при
таком изоморфизме нуль обязан переходить в элемент
$(0,0,0,\dots)$ (точнее, в класс этого элемента относительно
равенства), поскольку такой класс обладает свойствами нуля,
однозначно его определяющими (в~$\bbN$, а потому и
в~$\nstbb{\bbN}$). По аналогичным причинам единица переходит в
класс $(1,1,1,\dots)$ и вообще число $k$ соответствует классу
$(k,k,k,\dots)$. А класс $(0,1,2,3,4,\dots)$ отличается от любого
класса $(k,k,k,\dots)$ (они совпадают в единственном сомножителе,
а одноэлементное множество является малым, так как ультрафильтр
неглавный). Таким образом, построенная нами модель $\nstbb{\bbN}$
не является стандартной.

Аналогичное рассуждение позволяет построить и нестандартные
модели действительных чисел (о которых мы будем говорить в
следующем разделе).
